% Pregunta 901 (2003-julio-OEP2003, n.º 29): IS-LM-BP con tipo de cambio fijo y movilidad del capital
% limitada (BP más inclinada que la LM). Expansión fiscal y monetaria a la vez desde un punto de
% desempleo y déficit: la economía se aleja más de la BP y el déficit aumenta. Rectas a escala:
%   BP: i = 1,6 Y - 2,2   LM: i = 0,5 + 0,5 Y   IS: i = 5 - 0,9 Y
%   IS': i = 6,8 - 0,9 Y  LM': i = -0,1 + 0,5 Y
%   E0 = (3,21; 2,11), BP a ese tipo en Y = 2,69 (distancia 0,52)
%   E1 = (4,93; 2,36), BP a ese tipo en Y = 2,85 (distancia 2,08)   Ype = 5,55
% Autor: Víctor Gutiérrez Marcos
\documentclass[tikz,border=6pt]{standalone}
\input{tikz-tcee}
\begin{document}
\begin{tikzpicture}[tcee,x=1.15cm,y=1.15cm]
  % Zona de déficit (a la derecha y por debajo de la BP)
  \fill[area] (1.375,0) -- (4.375,4.8) -- (6,4.8) -- (6,0) -- cycle;
  \ejes{$Y$}{$i$}{6}{5}
  \node[etiqueta,font=\small\bfseries,text=tceeverde] at (5.0,4.55) {Déficit};
  \node[etiqueta,font=\small\bfseries,text=tceeverde] at (0.85,3.0) {Superávit};
  % Pleno empleo
  \draw[guia] (5.55,0) -- (5.55,3.75) node[above,text=tceegris] {$Y_{pe}$};
  % Distancia a la BP (tamaño del déficit), a la altura de cada equilibrio
  \draw[tceemorado,line width=2.2pt,opacity=0.55] (2.692,2.107) -- (3.214,2.107);
  \draw[tceerojo,line width=2.2pt,opacity=0.55] (2.852,2.364) -- (4.929,2.364);
  % Curvas iniciales
  \draw[curva3] (1.375,0) -- (4.375,4.8) node[etiqueta,above,text=tceeverde] {$BP$};
  \draw[curva1] (0.2,0.6) -- (5.6,3.3) node[etiqueta,right,text=tceemorado] {$LM$};
  \draw[curva2] (0.6,4.46) -- (5.4,0.14) node[etiqueta,right,text=tceeazul] {$IS$};
  % Lo que cambia
  \draw[curvanueva,desplazada] (2.2,4.82) -- (6.0,1.4) node[etiqueta,right,text=tceerojo] {$IS'$};
  \draw[curvanueva,desplazada] (0.4,0.1) -- (5.8,2.8) node[etiqueta,right,text=tceerojo] {$LM'$};
  % Proyecciones
  \draw[guia] (3.214,2.107) -- (3.214,0) node[below] {$Y_0$};
  \draw[guia] (4.929,2.364) -- (4.929,0) node[below] {$Y_1$};
  % Puntos
  \node[punto] at (3.214,2.107) {};
  \node[etiqueta,right=2pt] at (3.214,2.107) {$E_0$};
  \node[punto] at (4.929,2.364) {};
  \node[etiqueta,above=3pt] at (4.929,2.364) {$E_1$};
\end{tikzpicture}
\end{document}
